\section{Introduction}

% Starte generelt med å utforske hva den faglige diskusjonen går i for tiden
% Gå over til mer spesifikt, hva har andre gjort med problemstillingen som utforskes
% Presenter det vitenskapelige målet og problemstillingen, hva er det de andre ikke har fått til som vi skal gjøre

%generell del
Magnetic dipoles and domains have long served as the fundamental storage element in conventional data-storage devices, such as hard disks \cite{shukla2013_magnetic_moments_as_data_storage}. While these systems are robust at the macroscopic scale, the explosion in demand for data storage has driven device dimensions down to the nanoscale \cite{priyadarshi2023_devices_on_nanoscale}. At this scale, magnetic states become susceptible to thermal fluctuations, which can induce unwanted switching and limit stability \cite{everaert2023_mnp_noise}. 

In recent years, magnetic vortices have emerged as a promising alternative. A magnetic vortex is a spin structure where the magnetization rotates around a central core in a circular pattern \cite{behncke2018_magnetic_vortices}. In 2D, the in-plane magnetization circulates around the center, adopting one of two stable chiral states (clockwise (CW) or counterclockwise (CCW)) and, unlike conventional dipoles, these states do not switch spontaneously under thermal fluctuations \cite{sanchez‑manzano2024_stability_of_vortices}. This intrinsic stability has attracted significant research interest, positioning magnetic vortices as strong candidates for next-generation non-volatile memory cells in data-storage technologies \cite{spasojevic2025_vortices_promising_memory_cell}.

%spesifikk del
To fabricate such structures, several nanolithography techniques are available, including electron beam lithography (EBL), X-ray lithography (XRL), and focused ion beam (FIB). While EBL offers the highest resolution, and XRL is faster, both suffers from more complicated sample preparation \cite{ntnu_TMT4320_nanomat_lecture}. FIB has emerged as a versatile and increasingly popular alternative due to its high precision and ability to directly pattern materials without extensive sample preparation. One of its primary modes of operation is milling, where a focused Ga$^{+}$ ion beam is used to remove material by sputtering with nanometer-scale precision \cite{nist2021_fib_machining}, which is the mode used in this experiment.

%vitenskapelig mål/problemstilling
The objective of this paper is to investigate the effect of size and geometry of polygonal micro- and nanostructures on their magnetic domain configurations. More presicely, we will compare the average curl for figures of approximately 900, 600 and 400 nm in radii and look for a relationship as the number of sides increases, examining the transition from discrete to continuous magnetic domains in polygons as the number of sides vary. %Denne må endres til å passe resultater og diskusjon når vi har skrevet det.

%We will also discuss the figures' number of domains, and whether or not a vortex is formed, and if not, why this is the case. NEINEI?!